% ============================================================
%  ch09  第三讲 §3.2 分解性质、§3.3 Ran 空间、§3.4 $\Gr_{\Ran}$ 上的刚性化线丛
%  原文行 2014–2194
% ============================================================

\subsection{分解性质}

下面用定理 \ref{III:factorization} 概括 Beilinson--Drinfeld Grassmann 簇 $\{\Gr_{X^I}\}$ 最重要的性质，它是命题 \ref{III:Gr2} 的极大推广，但证明方式相似。

\begin{thm}\label{III:factorization}
设 $\underline G$ 是 $X$ 上光滑仿射群概形。

(i) 对每个 $I$，有单位截面 $e_I: X^I\to \Gr_{X^I}$，由平凡 $\underline G$-主丛给出。

(ii) 对每个满射 $\phi: J\twoheadrightarrow I$，以 $\Delta(\phi):X^I\to X^J$ 记相应的对角嵌入，它把 $\{x_i, i\in I\}$ 送到 $\{y_j, j\in J\}$，其中 $y_j=x_{\phi(j)}$。则有典范同构
\begin{equation}\label{BD restriction}
\Delta(\phi):\Gr_{X^I}\cong \Gr_{X^J}\times_{X^J,\Delta(\phi)}X^I.
\end{equation}
为简记，复合 $\Gr_{X^I}\cong \Gr_{X^J}\times_{X^J,\Delta(\phi)}X^I\to \Gr_{X^J}$ 也记作 $\Delta(\phi)$。

(iii) 给定满射 $\phi: J\twoheadrightarrow  I$，以 $J=\sqcup_{i\in I}J_i$ 记 $J$ 的分拆，以 $X^{(\phi)}\subset X^J$ 记由那些满足「$\phi(j)\neq \phi(j')$ 时 $x_{j}\cap x_{j'}=\emptyset$」的 $\{x_j, j\in J\}$ 组成的开子集。则有典范同构（通常称为分解同构）
\begin{equation}\label{BD factorization}
c_\phi: \Gr_{X^J}\times_{X^J}X^{(\phi)}\cong \Big(\prod_{i\in I}\Gr_{X^{J_i}}\Big)\times_{X^J}X^{(\phi)}.
\end{equation}
这些映射满足下述条件：
\begin{enumerate}
\item $\Delta(\phi)\circ e_I=e_J|_{\Delta(\phi)}$，且 $c_\phi (e_J|_{X^{(\phi)}})=(\prod_{i\in I} e_{J_i})|_{X^{(\phi)}}$。
\item 给定 $K\stackrel{\phi}{\twoheadrightarrow}  J\stackrel{\psi}{\twoheadrightarrow}  I$，有 $\Delta(\phi)\Delta(\psi)=\Delta(\psi\phi)$。
\item 给定 $K\stackrel{\phi}{\twoheadrightarrow}  J\stackrel{\psi}{\twoheadrightarrow}  I$，下图交换
\[\begin{CD}
\Gr_{X^J}|_{X^{(\psi)}}@>\Delta(\phi)|_{X^{(\psi)}}>>\Gr_{X^K}|_{X^{(\psi\phi)}}\\
@Vc_{\psi} VV@VV c_{\psi\phi}V\\
(\prod_i\Gr_{X^{J_i}})|_{X^{(\psi)}}@>(\prod_{i}\Delta(\phi_i))|_{X^{(\psi)}}>> (\prod_i \Gr_{X^{K_i}})|_{X^{(\psi\phi)}},
\end{CD}\]
其中 $\phi_i:K_i\to J_i$ 是 $i\in I$ 上的映射。
\item 给定 $K\stackrel{\phi}{\twoheadrightarrow}  J\stackrel{\psi}{\twoheadrightarrow}  I$，有 $(\prod_{i\in I}c_{\phi_i})(c_{\psi\phi}|_{X^{(\phi)}})=c_{\phi}$。
\end{enumerate}
\end{thm}

\begin{ex}\label{FA ex}
看看这个定理如何包含命题 \ref{III:Gr2} 的全部信息。首先，把 \eqref{BD restriction} 用于映射 $\phi:\{1,2\}\to \{1\}$，便得同构 $\Delta: \Gr_X\cong \Gr_{X^2}|_{\Delta}$；把它用于交换 $\sigma:\{1,2\}\to \{1,2\}$，便得覆盖 $X^2\to X^2$ 的交换映射的典范同构 $\Delta(\sigma):\Gr_{X^2}\to \Gr_{X^2}$。另一方面，把 \eqref{BD factorization} 用于映射 $\id: \{1,2\}\to \{1,2\}$，便得同构 $c:\Gr_{X^2}|_{X^2-\Delta}\cong \Gr_X\times\Gr_X|_{X^2-\Delta}$。把性质 (3) 用于 $\{1,2\}\stackrel{\sigma}{\to}\{1,2\}\stackrel{\id}{\to}\{1,2\}$、把性质 (4) 用于 $\{1,2\}\stackrel{\id}{\to}\{1,2\}\stackrel{\sigma}{\to}\{1,2\}$，可见同构 $c$ 与 $\Delta(\sigma)$ 以及 $\Gr_{X}\times\Gr_X$ 的自然交换相容，这正是命题 \ref{III:Gr2} 所断言者。
\end{ex}

命题 \ref{III:line} 的推广是

\begin{thm}\label{III:factor line}
设 $G$ 单连通，$\underline G=G\times X$ 常值。设 $\mL$ 是 $\Bun_{G}$ 上刚性化线丛，$\mL_{X^I}$ 是它到 $\Gr_{X^I}$ 的拉回。则在定理 \ref{III:factorization} 中的同构下，有典范同构

(i) $e_I^*\mL_{X^I}\cong \mO_{X^I}$；(ii) $\Delta(\phi)^*\mL_{X^J}\cong \mL_{X^I}$；(iii) $\mL_{X^J}|_{X^{(\phi)}}\cong c_\phi^*(\boxtimes \mL_{X^{J_i}}|_{X^{(\phi)}})$。

此外，这些同构彼此相容，即定理 \ref{III:factorization} 中的 (1)--(4) 有相应的陈述。
\end{thm}

上述 Beilinson--Drinfeld Grassmann 簇的结构与性质也为其他几何对象所共有，值得概括为下述定义。

\begin{dfn}\label{FA space}
$X$ 上的\textbf{分解空间} $\mY$ 是指：对每个非空有限集 $I$ 指定空间 $\mY_I\to X^I$，并给出一族如定理 \ref{III:factorization} (ii)--(iii) 的映射，满足该定理的相容条件 (2)--(4)。若还存在一族如定理 \ref{III:factorization} (i) 的映射并满足相容条件 (1)，则称该分解空间是\textbf{带单位的}。

带单位分解空间 $\mY$ 上的\textbf{线丛} $\mL$ 是指：对每个 $I$ 指定 $\mY_I$ 上的线丛 $\mL_I$，并给出一族如定理 \ref{III:factor line} (i)--(iii) 的同构，满足该定理中的自然相容条件。
\end{dfn}

\begin{rmk}
分解空间（相应地，带单位分解空间）在 \cite[\S\ 3.10.16]{BD0} 中称为手征半群（相应地，手征幺半群）。进一步讨论与例子见\emph{同上}。
\end{rmk}

\begin{rmk}\label{det factorization2}
也可以把命题 \ref{det factorization} 推广为：族 $\mL_{\det, X^I}$ 构成 $\Gr_{\GL_n, X^I}$ 上的分次分解线丛，它们之间有定理 \ref{III:factor line} (i)--(iii) 那样的一族同构，但这些同构之间的相容条件须作修改。更确切地说，定理 \ref{III:line} 中类似的陈述 (3) 须修改。
\end{rmk}

\subsection{Ran 空间}

把所有 $\{\Gr_{X^I}\}$ 组织成一个对象、并用它重述上述定理，在概念上更简洁。

\begin{dfn}
$X$ 的 \textbf{Ran 空间} $\Ran(X)$ 是这样的预层：它给每个 $R$ 配上 $X(R)$ 的有限非空子集构成的集合。
\end{dfn}

注意对任意有限非空集 $I$ 都有重言映射 $X^I\to \Ran(X)$，它诱导同构
\[\varinjlim_I X^I\to \Ran(X),\]
其中余极限取在所有有限非空集（态射为满射）的范畴上的预层范畴中。由于余极限所过的范畴并不滤过，$\Ran(X)$ 不是空间（虽然每个 $X^I$ 都是）。事实上，它连 \'etale 拓扑下的层都不是。

对 $x=\{x_i,i\in I\}\in \Ran(X)(R)$，以 $\Ga_x$ 记所有图形 $\Ga_{x_i}$ 之并（作为 $X_R$ 中的闭子概形）。

\begin{dfn}\label{dfn Ran Grass}
定义 \textbf{Ran Grassmann 簇} $\Gr_{\underline G,\Ran(X)}$ 为
\[\Gr_{\underline G,\Ran(X)}=\left\{(x,\mE,\beta)\ \left|\ \begin{array}{l}
x\in \Ran(X)(R),\quad \mE \text{ 是 } X_R \text{ 上的 } \underline{G}\text{-主丛},\\
\beta: \mE|_{X_R-\Ga_x}\cong \mE^0|_{X_R-\Ga_x} \text{ 是平凡化}
\end{array}\right.\right\}.\]
\end{dfn}

注意我们并不声称 $\Gr_{\underline G,\Ran(X)}$ 可表示，因为它\emph{不}可表示。不过有自然投影
\[q_{\Ran}: \Gr_{\underline G,\Ran(X)}\to \Ran(X)\]
与单位截面
\[e_{\Ran}:\Ran(X)\to  \Gr_{\underline G,\Ran(X)},\]
它们相对可表示。事实上
\[\Gr_{\underline G,X^I}=X^I\times_{\Ran(X),q_{\Ran}}\Gr_{\underline G,\Ran(X)}.\]
若曲线 $X$ 与群概形 $\underline G$ 在上下文中明确，为简洁把 $\Gr_{\underline G,\Ran(X)}$ 记作 $\Gr_{\Ran(X)}$ 或 $\Gr_{\Ran}$。

下面是编码 Beilinson--Drinfeld Grassmann 簇分解结构的简洁方式。首先，Ran 空间 $\Ran(X)$ 有由点集之并给出的半群结构
\[\on{union}: \Ran(X)\times \Ran(X)\to \Ran(X),\quad (\{x_i\},\{x_j\})\mapsto \{x_i,x_j\}.\]
以 $(\Ran\times \Ran)_{disj}\subset \Ran\times\Ran$ 记由满足 $\{x_i\}\cap \{x_j\}=\emptyset$ 的 $\{x_i\},\{x_j\}$ 组成的开子预层，类似地定义 $(\Ran\times \Ran\times \Ran)_{disj}$。于是定理 \ref{III:factorization} 可重述为

\begin{thm}\label{Ran Grass factorization}
有典范同构
\[a: \Gr_{\Ran}\times_{\Ran}(\Ran\times\Ran)_{disj}\simeq (\Gr_{\Ran}\times\Gr_{\Ran})|_{(\Ran\times\Ran)_{disj}},\]
它与单位截面 $e_{\Ran}:\Ran\to\Gr_{\Ran}$ 相容，并在 $(\Ran\times\Ran\times \Ran)_{disj}$ 上满足自然的余圈条件。
\end{thm}

验证定理 \ref{III:factorization} 与定理 \ref{Ran Grass factorization} 等价，留作练习。

\subsection{$\Gr_{\Ran}$ 上的刚性化线丛}\label{line bundle on BD Grass}

Ran Grassmann 簇 $\Gr_{\Ran}$ 是极引人入胜的几何对象。下一讲将说明它编码了 $X$ 上 $\underline G$-丛模空间的许多信息，因此研究 $\Gr_{\Ran}$ 的几何很重要。这里用一个例子说明方法：计算 $\Gr_{\underline{G}, \Ran(X)}$ 上刚性化线丛的 Picard 群胚，其中 $\underline{G}\to X$ 是 $X$ 上单连通单群概形。

对带截面 $e:S\to X$ 的预叠态射 $X\to S$，以 $\Pic^e(X/S)$ 记 $X$ 上沿线丛的 Picard 群胚，附带沿 $e$ 的平凡化。有下述定理。

\begin{thm}\label{Rigid line}
设 $G$ 单且单连通，$\underline G=G\times X$ 常值。则对任意点 $x\in X$，映射 $\Gr_x\to \Gr_{\Ran}$ 诱导同构 $\Pic^e(\Gr_{\Ran}/\Ran)\simeq \Pic^e(\Gr_x)$。
\end{thm}

注意
\[\Pic^e(\Gr_{\Ran}/\Ran)=\varinjlim\ \Pic^e(\Gr_{X^I}/X^I),\]
且 Picard 群胚 $\Pic^e(\Gr_{X^I}/X^I)$ 离散。此外它是 $(X^I)_{et}$ 上某个 \'etale 层的整体截面群
\[\mP ic^e(\Gr_{X^I}/X^I)(U)= \Pic^e((\Gr_{X^I})_U/U).\]
先研究 $\mP ic^e(\Gr_{X^I}/X^I)$。

\begin{lem}
设 $X$ 是光滑曲线，$\underline G$ 是 $X$ 上连通约化群概形，逐纤维连通且单连通。则有典范同构
\[\bZ_X\simeq \mP ic^e(\Gr_X/X),\]
其中 $\bZ_X$ 是 $X$ 上的常值层 $\bZ$。特别地，对每个点 $x\in X$，茎 $\mP ic^e(\Gr_X/X)_x=\Pic^e(\Gr_x)=\bZ$。
\end{lem}

\begin{proof}
先处理 $\underline G=G\times X$ 常值且存在 \'etale 态射 $X\to \bA^1$ 的情形。由于仿射 Grassmann 簇与 \'etale 基变换交换（见命题 \ref{III:property}），化归到 $X=\bA^1$。此时有典范同构 $\Gr_{\bA^1}\cong \Gr_0\times {\bA^1}$：因为 $\bA^1$ 上有整体坐标，$\Aut(D)$-主丛 $\widehat{\bA^1}\to \bA^1$ 平凡。引理随即成立。注意同构 $\bZ\cong \Pic^e(\Gr_{\bA^1})$ 是典范的，因为有丰沛生成元 $\mO(1)$。

一般地，存在 \'etale 映射 $f:U\to X$，使得 $\underline G\times_XU\cong G\times U$（对某个常值群 $G$），且存在 \'etale 映射 $U\to \bA^1$。同构 $\bZ_U\cong \mP ic^e(\Gr_U/U)$ 的典范性蕴含它下降为典范同构 $\bZ_X\to \mP ic^e(\Gr_X/X)$。
\end{proof}

在上述同构下，把 $n\in \bZ_X$ 对应的 $\Gr_X$ 上的线丛记作 $\mO(n)_X$。

关键一步是理解 $\mP ic^e(\Gr_{X^2}/X^2)$。注意由于 $\Gr_X\to X$ 归纳射影，由 See-Saw 原理（见命题 \ref{III:line} 的证明）有典范同构
\[ \mP ic^e(\Gr_X/X)\boxtimes \mP ic^e(\Gr_X/X)\cong\mP ic^e(\Gr_X\times\Gr_X/X\times X).\]
特别地
\[\mP ic^e(\Gr_{X^2}/X^2)|_{X^2-\Delta}\cong \mP ic^e(\Gr_X/X)\boxtimes \mP ic^e(\Gr_X/X)|_{X^2-\Delta}\cong(\bZ_X\boxtimes\bZ_X)|_{X^2-\Delta}.\]

\begin{lem}\label{Pic X2}
上述同构可延拓为正合列
\[0\to \mP ic^e(\Gr_{X^2}/X^2)\to \bZ_X\boxtimes\bZ_X\to \Delta_*\bZ_X\to 0, \]
其中映射 $\bZ_X\boxtimes\bZ_X\to \Delta_*\bZ_X$ 由 $(a,b)\mapsto a-b$ 给出。
\end{lem}

\begin{proof}
先证 $\mP ic^e(\Gr_{X^2}/X^2)\to  \bZ_X\boxtimes\bZ_X$ 单射，即它没有沿对角支撑的非零截面。设 $(x,x)$ 是对角上的点。如上，在 $X$ 上 \'etale 局部地可设 $\underline G=G\times X$ 常值，且 $\Gr_{X}\to X$ 同构于 $\Gr\times X\to X$。设 $V\to X\times X$ 是 $(x,x)$ 的 \'etale 开邻域，为简洁把 $V\cap\Delta$ 记作 $\Delta$。设 $\mL$ 是 $\Gr_V:=\Gr_{X^2}\times_{X^2}V$ 上的刚性化线丛，在 $\Gr_{V-\Delta}:=\Gr_V|_{V-\Delta}$ 上平凡。要证 $\mL$ 平凡，只需证它对每个 Schubert 簇 $\Gr_{\leq(\la,\mu)}|_V\subset \Gr_V$ 的限制平凡（定义见命题 \ref{III:closure}），为简洁记该限制为 $\mL$。

由刚性化线丛的定义，$\mL$ 带有平凡化 $\theta: \mO_{V}\cong e^*\mL$；又 $\mL|_{\Gr_{V-\Delta}}$ 平凡。由于沿映射 $\Gr_{\leq(\la,\mu)}|_{V-\Delta}\to V-\Delta$ 的拉回诱导正则函数空间的同构，存在唯一平凡化
\[s:  \mO_{\Gr_{\leq(\la,\mu)}|_{V-\Delta}}\cong \mL|_{V-\Delta},\]
使 $e^*(s)=\theta|_{V-\Delta}$。由命题 \ref{III:closure} 与定理 \ref{singularity}，$\Gr_{\leq(\la,\mu)}$ Cohen--Macaulay，故 $\mL((n-1)\Delta)$ 是 $\mL(n\Delta)$ 的子层。于是存在整数 $n$ 使
\[s\in \Gamma(\Gr_{\leq(\la,\mu)},\mL(n\Delta))\setminus \Gamma(\Gr_{\leq(\la,\mu)},\mL((n-1)\Delta)).\]
把 $s$ 看作 $\mL(n\Delta)$ 的正则截面。此时 $e^*(s)$ 与 $\mO_V\simeq e^*\mL\to e^*\mL(n\Delta)$ 作为 $(e^*\mL)(n\Delta)$ 的正则截面一致，特别地 $n\geq 0$。但若 $n>0$，则 $s$ 在 $e(\Delta)$ 处消失，从而必在整个 $\Gr_{\leq (\la,\mu)}|_\Delta$ 上消失（因为 $s$ 的零点集纯余维一且含于 $\Gr_{\leq (\la,\mu)}|_\Delta$，而后者由命题 \ref{III:closure} 仍是整概形）。这将蕴含 $s\in \Gamma(\Gr_{\leq(\la,\mu)},\mL((n-1)\Delta))$，矛盾。故 $n=0$。同理 $s$ 不沿 $e(\Delta)$ 消失，从而处处不消失，于是 $s$ 给出 $\mL$ 的平凡化。单射性证毕。

现设 $\mL$ 是 $\Gr_{X^2}$ 上的刚性化线丛，其在 $\Gr_{X^2}|_{X^2-\Delta}$ 上的限制为 $\mO(a)_X\boxtimes\mO(b)_X|_{X^2-\Delta}$。仍由 See-Saw 原理，它沿卷积映射 $m: \Gr_X\tilde\times\Gr_X\to \Gr_{X^2}$ 的拉回 $m^*\mL$ 形如 $\pr_1^*\mO(a)_X\otimes \mL'$，其中 $\mL'$ 是 $\Gr_X\tilde\times\Gr_X$ 上某个线丛，其在每个 $\pr_1^{-1}(y)(\simeq \Gr_X)$ 上的限制同构于 $\mO(b)_X$。在对角上的一个点处，整体卷积映射限制为局部卷积映射（见 \eqref{conv m0}）。于是简单计算 \eqref{conv m0} 之下 $\Gr$ 上线丛的拉回即可知 $a=b$。这表明复合 $\Pic^e(\Gr_{X^2}/X^2)\to \bZ_X\boxtimes\bZ_X\to \Delta_*\bZ_X$ 为零。

最后要证 $\Pic^e(\Gr_{X^2}/X^2)$ 满射到 $\ker(\bZ_X\boxtimes\bZ_X\to\Delta_*\bZ_X)$。同样只需考虑沿对角上的茎，故如前可过渡到 $(x,x)$ 的 \'etale 邻域，从而设 $X=\bA^1\subset \bP^1$，$\underline G=G\times X$ 常值。此时把 $\Bun_G(\bP^1)$ 上的线丛 $\mO(1)$（见引理 \ref{O1 on Bun}）拉回到 $\Gr_{X^2}$ 即可。
\end{proof}

同样的论证表明 $\mP ic^e(\Gr_{X^I}/X^I)$ 适合下述正合列
\[0\to \mP ic^e(\Gr_{X^I})\to \boxtimes_{i\in I}\bZ_X\to \bigoplus_{\phi:I\twoheadrightarrow J/\sim}\Delta(\phi)_*\boxtimes_{j\in J}\bZ_X.\]
这里 $|J|=|I|-1$，且 $\phi\sim \phi'$ 指存在同构 $J\cong J'$ 使 $\phi$ 与 $\phi'$ 相容。换言之，和取遍把 $I$ 分成 $|I|-1$ 个子集的所有分拆。映射 $d_{\phi}:\boxtimes_{i\in I}\bZ_X\to \Delta(\phi)_*\boxtimes_{j\in J}\bZ_X$ 如下给出：可写 $I=I'\cup\{1,2\}\to J=I'\cup\{1\}$，$d_\phi$ 在 $\boxtimes_{i\in I'}\bZ_X$-因子上为恒等，在对应于 $\{1,2\}\to\{1\}$ 的因子上为引理 \ref{Pic X2} 中的映射 $\bZ_X\boxtimes\bZ_X\to \Delta_*\bZ_X$。

\begin{cor}
对每个点 $x\in X$ 与每个满射 $\phi: J\twoheadrightarrow I$，映射 $\Pic^e(\Gr_X)\to \Pic^e(\Gr_x)$ 与 $\Pic^e(\Gr_{X^J}/X^J)\to\Pic^e(\Gr_{X^I}/X^I)$ 都是同构。
\end{cor}

定理 \ref{Rigid line} 随即成立。
